% figures/w01-6dof.svg — 6자유도: 힘·속도·위치가 어떻게 짝지어지는가.
%
%   bash _tools/tikz2svg.sh figures/src/w01-6dof.tex figures/w01-6dof.svg
%
% 3차원은 눈대중으로 기울이지 않는다 (standing-orders.md §3-2). TikZ 에 축 셋의
% 화면 사영을 직접 주면 좌표 (x,y,z) 를 그대로 쓸 수 있고, 사영은 TikZ 가 한다.
% 사영은 _tools/w01_6dof_arcs.awk 가 쓰던 것과 같다:
%
%   P(x) = ( 0.87687, -0.21477)   선수  — 오른쪽 위 (화면 아래 성분이 음수)
%   P(y) = ( 0.48072,  0.39176)   우현  — 오른쪽 아래
%   P(z) = ( 0.00000,  0.89465)   아래
%
% TikZ 의 y 는 위가 양수이므로 두 번째 성분의 부호를 뒤집어 넣는다.
%
% 모멘트 호는 그 축에 수직인 평면 위의 원이다. 3차원 좌표로 원을 매개하면
% 사영된 타원과 그 중심이 축선 위에 오는 것이 저절로 보장된다 — 이것이
% 손으로 타원을 놓지 않는 이유다.
\documentclass[border=5pt]{standalone}
% gnc-style.tex — 이 볼트의 TikZ 그림이 공유하는 스타일.
%
% 저널 그림 규약을 스타일 하나로 굳혀 둔다. 그림마다 선 굵기나 화살촉을
% 다시 정하지 않는다 — 그것이 그림이 제각각으로 보이는 원인이다.
%
% 쓰는 법:  % gnc-style.tex — 이 볼트의 TikZ 그림이 공유하는 스타일.
%
% 저널 그림 규약을 스타일 하나로 굳혀 둔다. 그림마다 선 굵기나 화살촉을
% 다시 정하지 않는다 — 그것이 그림이 제각각으로 보이는 원인이다.
%
% 쓰는 법:  % gnc-style.tex — 이 볼트의 TikZ 그림이 공유하는 스타일.
%
% 저널 그림 규약을 스타일 하나로 굳혀 둔다. 그림마다 선 굵기나 화살촉을
% 다시 정하지 않는다 — 그것이 그림이 제각각으로 보이는 원인이다.
%
% 쓰는 법:  \input{gnc-style}  (figures/src/ 안에서)

\usepackage{tikz}
\usepackage{amsmath}
\usepackage{xcolor}
\usetikzlibrary{arrows.meta, positioning, calc, fit, backgrounds, decorations.pathmorphing}

% ---- 색 --------------------------------------------------------------------
% 색은 강조 하나에만 쓴다. 나머지는 검정.
\definecolor{ink}{HTML}{1A1A1A}
\definecolor{accent}{HTML}{7C3AED}
\definecolor{warn}{HTML}{B91C1C}
\definecolor{soft}{HTML}{666666}

% ---- 선과 화살촉 -----------------------------------------------------------
\tikzset{
  % 신호선. 화살촉은 작고 뾰족한 것 하나뿐이다.
  sig/.style   = {ink, line width=0.5pt, -{Stealth[length=4pt, width=3pt]}},
  wire/.style  = {ink, line width=0.5pt},
  % 강조 경로 — 파선으로 구분한다. 흑백 인쇄에서도 살아야 하기 때문이다.
  asig/.style  = {accent, line width=0.5pt, densely dashed,
                  -{Stealth[length=4pt, width=3pt]}},
  awire/.style = {accent, line width=0.5pt, densely dashed},
  % 블록. 흰 바탕, 직각, 검은 테두리.
  blk/.style   = {draw=ink, line width=0.55pt, fill=white,
                  minimum width=17mm, minimum height=8mm, inner sep=2pt, font=\small},
  % 합산점
  sum/.style   = {draw=ink, line width=0.55pt, fill=white, circle,
                  inner sep=0pt, minimum size=4.6mm, font=\scriptsize},
  asum/.style  = {draw=accent, line width=0.55pt, fill=white, circle,
                  inner sep=0pt, minimum size=4.6mm, font=\scriptsize},
  % 분기점
  dot/.style   = {ink, fill=ink, circle, inner sep=0pt, minimum size=1.5mm},
  % 신호 이름 — 선 위에 작게
  sname/.style = {font=\small, inner sep=1.5pt},
  % 그림 안의 짧은 설명. 그림 안의 글은 최소로 둔다
  note/.style  = {font=\scriptsize, soft, inner sep=1.5pt},
  anote/.style = {font=\scriptsize, accent, inner sep=1.5pt},
  % 축
  axis/.style  = {ink, line width=0.5pt},
  tick/.style  = {font=\scriptsize, soft},
}

% 캡션 한 줄 — 그림 아래 가는 선과 함께
\newcommand{\gnccaption}[2]{%
  \node[anchor=north west, text width=#1, align=left, font=\scriptsize, ink]
        at (current bounding box.south west) {\vspace{2pt}#2};}
  (figures/src/ 안에서)

\usepackage{tikz}
\usepackage{amsmath}
\usepackage{xcolor}
\usetikzlibrary{arrows.meta, positioning, calc, fit, backgrounds, decorations.pathmorphing}

% ---- 색 --------------------------------------------------------------------
% 색은 강조 하나에만 쓴다. 나머지는 검정.
\definecolor{ink}{HTML}{1A1A1A}
\definecolor{accent}{HTML}{7C3AED}
\definecolor{warn}{HTML}{B91C1C}
\definecolor{soft}{HTML}{666666}

% ---- 선과 화살촉 -----------------------------------------------------------
\tikzset{
  % 신호선. 화살촉은 작고 뾰족한 것 하나뿐이다.
  sig/.style   = {ink, line width=0.5pt, -{Stealth[length=4pt, width=3pt]}},
  wire/.style  = {ink, line width=0.5pt},
  % 강조 경로 — 파선으로 구분한다. 흑백 인쇄에서도 살아야 하기 때문이다.
  asig/.style  = {accent, line width=0.5pt, densely dashed,
                  -{Stealth[length=4pt, width=3pt]}},
  awire/.style = {accent, line width=0.5pt, densely dashed},
  % 블록. 흰 바탕, 직각, 검은 테두리.
  blk/.style   = {draw=ink, line width=0.55pt, fill=white,
                  minimum width=17mm, minimum height=8mm, inner sep=2pt, font=\small},
  % 합산점
  sum/.style   = {draw=ink, line width=0.55pt, fill=white, circle,
                  inner sep=0pt, minimum size=4.6mm, font=\scriptsize},
  asum/.style  = {draw=accent, line width=0.55pt, fill=white, circle,
                  inner sep=0pt, minimum size=4.6mm, font=\scriptsize},
  % 분기점
  dot/.style   = {ink, fill=ink, circle, inner sep=0pt, minimum size=1.5mm},
  % 신호 이름 — 선 위에 작게
  sname/.style = {font=\small, inner sep=1.5pt},
  % 그림 안의 짧은 설명. 그림 안의 글은 최소로 둔다
  note/.style  = {font=\scriptsize, soft, inner sep=1.5pt},
  anote/.style = {font=\scriptsize, accent, inner sep=1.5pt},
  % 축
  axis/.style  = {ink, line width=0.5pt},
  tick/.style  = {font=\scriptsize, soft},
}

% 캡션 한 줄 — 그림 아래 가는 선과 함께
\newcommand{\gnccaption}[2]{%
  \node[anchor=north west, text width=#1, align=left, font=\scriptsize, ink]
        at (current bounding box.south west) {\vspace{2pt}#2};}
  (figures/src/ 안에서)

\usepackage{tikz}
\usepackage{amsmath}
\usepackage{xcolor}
\usetikzlibrary{arrows.meta, positioning, calc, fit, backgrounds, decorations.pathmorphing}

% ---- 색 --------------------------------------------------------------------
% 색은 강조 하나에만 쓴다. 나머지는 검정.
\definecolor{ink}{HTML}{1A1A1A}
\definecolor{accent}{HTML}{7C3AED}
\definecolor{warn}{HTML}{B91C1C}
\definecolor{soft}{HTML}{666666}

% ---- 선과 화살촉 -----------------------------------------------------------
\tikzset{
  % 신호선. 화살촉은 작고 뾰족한 것 하나뿐이다.
  sig/.style   = {ink, line width=0.5pt, -{Stealth[length=4pt, width=3pt]}},
  wire/.style  = {ink, line width=0.5pt},
  % 강조 경로 — 파선으로 구분한다. 흑백 인쇄에서도 살아야 하기 때문이다.
  asig/.style  = {accent, line width=0.5pt, densely dashed,
                  -{Stealth[length=4pt, width=3pt]}},
  awire/.style = {accent, line width=0.5pt, densely dashed},
  % 블록. 흰 바탕, 직각, 검은 테두리.
  blk/.style   = {draw=ink, line width=0.55pt, fill=white,
                  minimum width=17mm, minimum height=8mm, inner sep=2pt, font=\small},
  % 합산점
  sum/.style   = {draw=ink, line width=0.55pt, fill=white, circle,
                  inner sep=0pt, minimum size=4.6mm, font=\scriptsize},
  asum/.style  = {draw=accent, line width=0.55pt, fill=white, circle,
                  inner sep=0pt, minimum size=4.6mm, font=\scriptsize},
  % 분기점
  dot/.style   = {ink, fill=ink, circle, inner sep=0pt, minimum size=1.5mm},
  % 신호 이름 — 선 위에 작게
  sname/.style = {font=\small, inner sep=1.5pt},
  % 그림 안의 짧은 설명. 그림 안의 글은 최소로 둔다
  note/.style  = {font=\scriptsize, soft, inner sep=1.5pt},
  anote/.style = {font=\scriptsize, accent, inner sep=1.5pt},
  % 축
  axis/.style  = {ink, line width=0.5pt},
  tick/.style  = {font=\scriptsize, soft},
}

% 캡션 한 줄 — 그림 아래 가는 선과 함께
\newcommand{\gnccaption}[2]{%
  \node[anchor=north west, text width=#1, align=left, font=\scriptsize, ink]
        at (current bounding box.south west) {\vspace{2pt}#2};}


\begin{document}
\begin{tikzpicture}[
    x={(0.87687cm,0.21477cm)},
    y={(0.48072cm,-0.39176cm)},
    z={(0cm,-0.89465cm)},
    scale=3.4]

% ---- 축 --------------------------------------------------------------------
\draw[axis, -{Stealth[length=4.5pt,width=3.4pt]}] (0,0,0) -- (1.5,0,0);
\draw[axis, -{Stealth[length=4.5pt,width=3.4pt]}] (0,0,0) -- (0,1.5,0);
\draw[axis, -{Stealth[length=4.5pt,width=3.4pt]}] (0,0,0) -- (0,0,1.5);
\node[font=\scriptsize, soft] at (1.72,0,0)  {$x_b$};
\node[font=\scriptsize, soft] at (0,1.66,0)  {$y_b$};
\node[font=\scriptsize, soft] at (0,0,1.66)  {$z_b$};
\fill[ink] (0,0,0) circle (0.6pt);
\node[font=\scriptsize, soft, anchor=east] at (-0.12,0,0) {$o_b$};

% ---- 갑판 (선수가 +x) -------------------------------------------------------
\fill[ink, opacity=0.10]
      (0.62,0,0) -- (0.34,0.26,0) -- (-0.53,0.26,0) -- (-0.62,0.15,0)
   -- (-0.62,-0.15,0) -- (-0.53,-0.26,0) -- (0.34,-0.26,0) -- cycle;
\draw[soft, line width=0.3pt]
      (0.62,0,0) -- (0.34,0.26,0) -- (-0.53,0.26,0) -- (-0.62,0.15,0)
   -- (-0.62,-0.15,0) -- (-0.53,-0.26,0) -- (0.34,-0.26,0) -- cycle;

% ---- 힘 셋 (직선 화살표) ----------------------------------------------------
\draw[alng, line width=1.0pt, -{Stealth[length=5pt,width=3.6pt]}]
      (0.75,0,0) -- (1.28,0,0);
\node[alng, font=\scriptsize] at (1.14,0.22,0) {$X$ — surge};
\draw[alng, line width=1.0pt, -{Stealth[length=5pt,width=3.6pt]}]
      (0,0.45,0) -- (0,1.28,0);
\node[alng, font=\scriptsize, anchor=north] at (0,1.20,0.22) {$Y$ — sway};
\draw[alng, line width=1.0pt, -{Stealth[length=5pt,width=3.6pt]}]
      (0,0,0.30) -- (0,0,1.28);
\node[alng, font=\scriptsize, anchor=west] at (0.10,0,1.24) {$Z$ — heave};

% ---- 모멘트 셋 (축에 수직인 평면 위의 원) ----------------------------------
% x 축 둘레 : y -> z 방향으로 도는 원, 중심은 x 축 위
\draw[accent, line width=0.9pt, -{Stealth[length=4.5pt,width=3.4pt]}]
      plot[domain=40:340, samples=60, variable=\t]
      ({2.05}, {0.30*cos(\t)}, {0.30*sin(\t)});
\node[accent, font=\scriptsize, anchor=south] at (2.05,0,-0.42) {$K$ — roll};

% y 축 둘레 : z -> x
\draw[accent, line width=0.9pt, -{Stealth[length=4.5pt,width=3.4pt]}]
      plot[domain=40:340, samples=60, variable=\t]
      ({0.30*sin(\t)}, {1.85}, {0.30*cos(\t)});
\node[accent, font=\scriptsize, anchor=south] at (0,1.85,-0.42) {$M$ — pitch};

% z 축 둘레 : x -> y
\draw[accent, line width=0.9pt, -{Stealth[length=4.5pt,width=3.4pt]}]
      plot[domain=40:340, samples=60, variable=\t]
      ({0.30*cos(\t)}, {0.30*sin(\t)}, {1.90});
\node[accent, font=\scriptsize, anchor=west] at (0,0.52,1.90) {$N$ — yaw};

% 축선을 호 중심까지 점선으로 — "이 축 둘레" 임을 보이게 한다
\draw[soft, line width=0.3pt, densely dotted] (1.5,0,0) -- (2.05,0,0);
\draw[soft, line width=0.3pt, densely dotted] (0,1.5,0) -- (0,1.85,0);
\draw[soft, line width=0.3pt, densely dotted] (0,0,1.5) -- (0,0,1.90);
\fill[accent] (2.05,0,0) circle (0.4pt);
\fill[accent] (0,1.85,0) circle (0.4pt);
\fill[accent] (0,0,1.90) circle (0.4pt);

\end{tikzpicture}

\hspace{4mm}
% ===================== 오른쪽 : 표 =========================================
\begin{tikzpicture}[x=1mm, y=1mm]
\node[ink, font=\small\bfseries, anchor=north west] at (0,64)
     {Six degrees of freedom};
\node[anchor=north west, align=left, text width=76mm, font=\scriptsize, soft] at (0,58) {%
  Each row is one direction the vessel can move. Read it across: a force produces
  a velocity, and integrating that velocity gives a position.};

\node[anchor=north west, font=\scriptsize] at (0,44) {%
  \begin{tabular}{@{}llll@{}}
    \textbf{DOF} & \textbf{force / moment} & \textbf{lin. / ang. vel.} & \textbf{position / angle}\\[2pt]
    1\ \ surge & $X$ & $u$ & $x$\\
    2\ \ sway  & $Y$ & $v$ & $y$\\
    3\ \ heave & $Z$ & $w$ & $z$\\
    4\ \ roll  & $K$ & $p$ & $\phi$\\
    5\ \ pitch & $M$ & $q$ & $\theta$\\
    6\ \ yaw   & $N$ & $r$ & $\psi$\\
  \end{tabular}};

\node[anchor=north west, align=left, text width=76mm, font=\scriptsize] at (0,-4) {%
  \textbf{Straight arrow} = force. \textbf{Curved arrow} = moment, drawn as a
  circle in the plane perpendicular to its own axis, by the right-hand rule.
  The dotted stub joins each circle's centre to its own axis — an ellipse whose
  centre is off the axis is the classic way this picture goes wrong.\\[4pt]
  A surface craft is normally reduced to DOF 1, 2 and 6: surge, sway and yaw.
  Heave, roll and pitch remain in \texttt{otter.m} and appear in §1-12 as the
  reason $M_{66}$ is not simply $I_z$.};
\end{tikzpicture}

\end{document}
